\chapter{Introduction}
\label{cha:Introduction}
%\lipsum \autocite{DummyCite}

Large Language Models (LLMs) have fundamentally transformed how we look up information and obtain answers to questions, becoming an increasingly important tool, not only in everyday life, but also across a range of professional fields. With the usage of such systems, the need to verify the factual correctness and easily do further research on the underlying sources arises, especially when answers lack important details or provide a depth of explanation that is misaligned with the user's existing knowledge. This is especially important in high-stakes fields such as medicine, law or scientific research, where seemingly plausible but unbacked or false claims could lead users to draw false conclusions.

One increasingly popular method is the approach of Retrieval Augmented Generation (RAG), replacing zero-shot LLM solutions with a two-step process. Firstly, relevant sources in the form of documents or other multi-media data are retrieved and then fed into the generating model for it to use as ground truth for answer generation and later to enable citation extraction. 
Retrieving correct and precise text-spans alongside the generated answer is not a trivial task, and with domain-specific requirements widely varying, such systems are part of active research \cite{aSurveyOfEvidenceBasedTextGenerationWithLargeLanguageModels} \cite{citationAKeyToBuildingResponsibleAndAccountableLargeLanguageModels}.
One system applying such a RAG approach is SQuAI \cite{SQuAI}, aimed at answering complex open-domain scientific questions on its knowledgebase containing 2.3 million full-text scientific papers. For further performance improvements, it introduces a multi-agent pipeline, enhancing document retrieval capabilities as well as providing document-level backlinks for parts of its generated answer. This system, while providing some backlinking, does still fall short of providing more granular, sentence-level cited text spans, which have proven to be much more useful to users, as they are  quicker and easier to verify, enhancing trustworthiness and usability of the system.

\begin{center}    
    \includegraphics[page=1, width=1\textwidth]{citedTextSpanExtraction.pdf} 
    \captionof{figure}{Sentence level cited text span attribution}
\end{center}

Therefore, this thesis will focus on benchmarking different cited text span extraction methods to enhance source attribution, leading to the following core contributions:
\begin{enumerate}
    \item Curate a dataset of high-quality, multi-hop scientific questions building upon a large-scale retrieval corpus that require complex reasoning capabilities in order to be answered and closely resemble real-world scientific inquiry needs.
    \item Introduce a hierarchy of different cited text span extraction methods ranging from simple lexical and semantic matching up to complex reasoning capabilities, as well as strategic combinations of such, to allow for tradeoffs in accuracy and efficiency.
    \item Apply an automated framework to benchmark the retrieval and answering capabilities of non-parametric, retrieval- and evidence-based text generation systems on the proposed dataset as well as evaluate the extraction methods in terms of efficiency, correctness and quality of source attribution using NLI models and LLM-as-a-judge approaches.
\end{enumerate}




%noch verbessern:
%mmmh naja: The RAG approach itself however, does still fall short of providing granular proof for the models answer, at best providing a source at document-level.   
% maybe add illustration for squai
% mention unarXive 2022
% depending on if human evaluation plays a part add that


%additions
% PeerQA paper: number of scientific articles is increasing exponentially --> need new tools to stay up to date


%
% hallucination
% consequences
% need for verification 
% LLM good at synthesizing, but prone to hallucination
%
% hard to verify because of lacking domain knowledge, time, no knowledge about sources false conclusions without any data
%RAG - briefly mention squai
% - does not link back to specific text spans
%briefly mention post vs pre hoc (squai post) - text based
% what im planning on solving, we cant even measure without a good dataset / expensive human evaluation --> curate that
% how to improve & still get good performance (keyword, semantic --> increasing complexity while still having ), tradeof
%   -dataset - extaction -evalutino
% maybe example
% 
%


%Things to mention in the thesis:
%alternatives to rag: -iterative feedback loops -supervised fine-tuning -prompt engineering -self-reflection methodologies
%
%
%
%


% models are nutoriously bad at multitasking (pre hoc, take as exmaple dataset citation generation)
% lexical missmatch --> could not realise ambiguity